\documentclass[12pt]{article}

\usepackage[spanish]{babel}
\usepackage[utf8]{inputenc}
\usepackage[T1]{fontenc}
\usepackage{amsmath, amssymb, bm}
\usepackage{geometry}
\usepackage{setspace}
\usepackage{booktabs}

\geometry{margin=2.5cm}
\onehalfspacing

\title{Interpretación de Resultados en el Modelo SFA}
\author{}
\date{}

\begin{document}

\maketitle

\section{El punto de partida: qué mide $u_{it}$ y por qué necesitamos dos}

Antes de entrar en los contrastes específicos, hay que entender qué estás estimando realmente. En el modelo SFA, $u_{it}$ es la distancia entre lo que el hospital produce y lo que podría producir con los mismos inputs si operara en la frontera. No es un residuo estadístico cualquiera: es la parte del error que, por construcción del modelo, solo puede ser positiva y que se interpreta como ineficiencia técnica.

La clave de toda la estrategia empírica es que el modelo teórico predice efectos asimétricos: la descentralización debería reducir la ineficiencia en cantidad y aumentarla en intensidad. Si se utilizara un único $u_{it}$ para un output agregado, estos dos efectos se cancelarían algebraicamente. Por ello, es necesario estimar dos fronteras con dos términos $u_{it}$ diferenciados: uno para cantidad y otro para intensidad.

\section{La ecuación de ineficiencia}

En cada frontera, la media del término de ineficiencia es:

\begin{equation}
\mu_{kit} =
\underbrace{\delta_0}_{\text{nivel base}} +
\underbrace{\delta_1 D^{desc}_{it}}_{\text{estructura org.}} +
\underbrace{\delta_2 pct^{SNS}_{it}}_{\text{mecanismo pago}} +
\underbrace{\delta_3 (D^{desc} \times pct^{SNS})_{it}}_{\text{interacción}} +
\underbrace{\delta_4 ShareQ_{it}}_{\text{proxy } \psi_{12}} +
\underbrace{\delta_5 (D^{desc} \times ShareQ)_{it}}_{\text{moderación}} +
\text{CCAA} + \text{cluster}
\end{equation}

Un $\mu_{kit}$ más negativo implica menor ineficiencia esperada. Los coeficientes miden el efecto de cada variable sobre la media de la distribución de ineficiencia.

Al estimar dos fronteras, se obtienen dos vectores de coeficientes: $\bm{\alpha}$ (cantidad o quirúrgico) y $\bm{\delta}$ (intensidad o médico). La cuestión central es si $\bm{\alpha} \neq \bm{\delta}$.

\section{Predicción P1: descentralización y cantidad}

El modelo predice que la descentralización reduce la ineficiencia en cantidad.

\textbf{Test:}
\[
H_0: \alpha_1 = 0 \quad \text{vs} \quad H_1: \alpha_1 < 0
\]

\textbf{Resultados:}
\begin{itemize}
\item Diseño A: $\hat{\alpha}_1 = -2.530$, $t = -6.98$
\item Diseño B: $\hat{\alpha}_{Merc} = +2.256$, $t = +5.08$
\end{itemize}

El resultado del Diseño A implica una reducción de la ineficiencia de aproximadamente 4--6 puntos porcentuales.

El signo positivo en el Diseño B no contradice la hipótesis: refleja especialización (ineficiencia en cirugía, eficiencia en medicina).

\textbf{Conclusión:} P1 confirmada robustamente.

\section{Predicción P2: descentralización e intensidad}

El modelo predice que la descentralización aumenta la ineficiencia en intensidad.

\textbf{Test:}
\[
H_0: \delta_1 = 0 \quad \text{vs} \quad H_1: \delta_1 > 0
\]

\textbf{Resultado:}
\[
\hat{\delta}_1 = -1.477, \quad t = -4.21
\]

No se confirma el signo esperado.

Sin embargo, el contraste relevante es:

\[
H_0: \alpha_1 = \delta_1 \quad \text{vs} \quad H_1: \alpha_1 < \delta_1
\]

\[
z_{dif} = -2.09, \quad p = 0.037
\]

\textbf{Conclusión:} evidencia parcial. Existe asimetría aunque no inversión de signo.

\section{Predicción P3: efecto del tipo de pago}

\textbf{Resultados:}
\begin{itemize}
\item $pct^{SNS}$: $z_{dif} = -0.088$, $p = 0.930$
\item Interacción: $z_{dif} = +2.349$, $p = 0.019$
\end{itemize}

El efecto diferencial se encuentra en la interacción, no en el coeficiente principal.

En el Diseño B:
\[
z_{dif} = +5.32, \quad p < 0.001
\]

\textbf{Conclusión:} P3 confirmada.

\section{Predicción P4: papel de $\psi_{12}$}

\textbf{Resultados:}

\begin{itemize}
\item Diseño A: no significativo
\item Diseño B: signos opuestos y altamente significativos
\end{itemize}

\[
\hat{\alpha}_4 = -7.015, \quad \hat{\delta}_4 = +4.647
\]

\textbf{Conclusión:} P4 confirmada en Diseño B.

\section{Predicción P*: contraste global}

\textbf{Test:}
\[
LR = 2(\ell^{irr} - \ell^{restr}) \sim \chi^2(k)
\]

\textbf{Resultados:}
\begin{itemize}
\item Diseño A: $LR = 10490.8$
\item Diseño B: $LR = 2508.0$
\end{itemize}

\textbf{Conclusión:} rechazo contundente de la igualdad de coeficientes.

\section{Correlación entre ineficiencias}

\[
r = -0.140, \quad p < 0.001
\]

Interpretación: existe sustitución de esfuerzo entre dimensiones.

\section{Resumen de resultados}

\begin{center}
\begin{tabular}{lllll}
\toprule
Predicción & Mecanismo & Diseño A & Diseño B & Veredicto \\
\midrule
P1 & $\downarrow$ ineficiencia en cantidad & Confirmada & Confirmada & Sí \\
P2 & $\uparrow$ ineficiencia en intensidad & Parcial & Parcial & Parcial \\
P3 & Efectos asimétricos pago & Parcial & Fuerte & Sí \\
P4 & Moderación $\psi_{12}$ & No & Sí & Sí (B) \\
P* & Vectores distintos & Sí & Sí & Sí \\
\bottomrule
\end{tabular}
\end{center}

\end{document}